\section{Results and Discussion}
\label{sec:results_discussion}

The campaigns follow the protocol of Section~\ref{subsec:numerical_protocol} and are read against the many-query regime of Section~\ref{sec:introduction}: cost is reported per forward solve, because that is the quantity such a calculation multiplies, whereas the factorisation of $\mathbf{A}$ is performed once and amortised over the whole calculation. Four questions are addressed. Does the local spectral radius identify the regime in which Picard reaches the manufactured reference equilibrium? How does the nonlinear work change as $\rho=1$ is approached? At which state can the spectral radius be evaluated and still carry that information? And does updating it during the nonlinear iteration repay its cost?

Four campaigns are reported, one strategy at a time and over the same spectral partition; comparisons between strategies are deferred to Section~\ref{subsec:many_query}, where the paired design places them on a common basis. Each cell gives three quantities: $C_{\rm rel}$, the median algorithmic time per case normalised by the smallest median within that campaign; $\overline{N}_{\rm lin}$, the mean number of linear systems solved by the nonlinear method, rounded to the nearest integer; and $P_{\rm ref\text{-}fail}$, the observed probability of failing to reach the manufactured reference solution among cases for which $\rho_{\rm ref}$ is available, which for Picard covers both residual non-convergence and convergence to a residual-satisfying state distinct from that solution. The two mixed campaigns add the fraction of nonlinear iterations performed by Picard. The technical exceptions of Section~\ref{subsec:numerical_protocol} carry no $\rho_{\rm ref}$ and are assigned to no spectral bin. The normalising time of each campaign is the smallest of its $12$ cell medians: $0.125103$, $0.067639$, $0.139202$, and $0.275310\,\mathrm{s}$ for the Picard-only, Newton-only, initial-$\rho$, and updated-$\rho$ campaigns respectively. Section~\ref{subsec:many_query} instead normalises by a median over all cases of one cell, and the two bases are not interchangeable. The results below are for WLSQR.

\begin{table}[htb]
\centering
\caption{How Picard's cost, work and failure rate change across the spectral partition.}
\label{tab:final_picard}
\setlength{\tabcolsep}{7pt}
\begin{tabular}{lcc}
\toprule
\rowcolor{gray!12}
Mesh &
$\rho<0.90$ &
$0.90\leq\rho<0.98$ \\
\midrule
Cartesian &
$(1.00,\,29,\,0.0\%)$ &
$(2.90,\,83,\,0.0\%)$ \\
Hexagonal &
$(1.17,\,29,\,0.0\%)$ &
$(3.24,\,81,\,0.0\%)$ \\
Random &
$(1.44,\,30,\,0.0\%)$ &
$(3.98,\,86,\,0.0\%)$ \\

\addlinespace[2pt]
\rowcolor{gray!12}
Mesh &
$0.98\leq\rho<1.02$ &
$\rho\geq1.02$ \\
\midrule
Cartesian &
$(21.52,\,1205,\,45.6\%)$ &
$(8.15,\,295,\,100.0\%)$ \\
Hexagonal &
$(23.82,\,1185,\,45.7\%)$ &
$(9.23,\,292,\,100.0\%)$ \\
Random &
$(28.28,\,1075,\,40.2\%)$ &
$(10.51,\,290,\,100.0\%)$ \\
\bottomrule
\end{tabular}
\end{table}




\begin{table}[htb]
\centering
\caption{The outcome of each Picard solve, and the size of the error when it is wrong.}
\label{tab:branch_classification}
\setlength{\tabcolsep}{7pt}
\begin{tabular}{lcccc}
\toprule
\rowcolor{gray!12}
 & & Reached & Other & Did not \\
\rowcolor{gray!12}
 & Cases & reference & state & converge \\
\midrule
$\rho_{\rm ref}<1$ & 3019 & 3019 & 0 & 0 \\
$\rho_{\rm ref}\geq1$ & 2737 & 46 & 1696 & 995 \\
\addlinespace[4pt]
\rowcolor{gray!12}
 & & Median & Maximum & \\
\rowcolor{gray!12}
$\rho_{\rm ref}$ & Cases & distance & distance & \\
\midrule
$1.00$--$1.05$ & 541 & $0.40\%$ & $1.11\%$ & \\
$1.05$--$1.10$ & 418 & $1.46\%$ & $2.29\%$ & \\
$1.10$--$1.20$ & 338 & $2.94\%$ & $4.49\%$ & \\
$1.20$--$1.40$ & 396 & $4.82\%$ & $6.68\%$ & \\
$1.40$--$2.35$ & 3 & $0.19\%$ & $0.39\%$ & \\
\bottomrule
\end{tabular}
\end{table}

\subsection{The spectral transition of the fixed-point iteration}
\label{subsec:picard_transition}

Table~\ref{tab:final_picard} shows the transition of the fixed-point iteration. The subcritical regime remains robust, whereas the nonlinear work rises by a factor of between $36$ and $42$ between the lowest bin and the critical interval. Once the reference equilibrium becomes locally non-attractive, failures to reach it appear and then dominate. Cost falls again above $\rho=1.02$ because the bin mixes two terminations: $758$ of the $2188$ supercritical cases reach the divergence threshold and stop early, while the remaining $1430$ converge, to the reference or elsewhere. In the critical interval almost every case converges, and slowly. The same transition is observed on all three mesh families, which places it with the spectral regime rather than with a particular mesh topology.

\begin{figure}[h]
  \centering
  \includegraphics[width=0.45\textwidth]{Figures/Figure_01.png}
\caption{Distance $d_{\rm ref}$ of the alternative solutions from the reference equilibrium against $\rho_{\rm ref}-1$, by mesh family.}
  \label{fig:branch_scaling}
\end{figure}

\paragraph{Residual convergence to a different discrete solution}
Most of these failures are of one kind. Of the $2691$ Picard solves that did not reach the reference equilibrium, $1696$, or $63.0\%$, satisfied the nonlinear residual criterion at a different state. Table~\ref{tab:branch_classification} sets the outcome against $\rho_{\rm ref}$, the spectral radius at the reference equilibrium, over the $5756$ of $6000$ WLSQR cases for which it is available. Below unity there is no exception. Above it the cases divide between a different discrete state and outright non-convergence, with a remainder that still reaches the reference; that remainder lies entirely below $\rho_{\rm ref}=1.015544$, and the smallest $\rho_{\rm ref}$ at which any case failed is $1.000290$, so the whole transition is confined to an interval of width $0.0153$.


The non-reference states are solutions of the discrete system and not premature terminations: their final residuals lie between $0.77$ and $1.01$ times $10^{-10}$, within the test of Table~\ref{tab:protocol}, whose right-hand side is at least $1.01\times10^{-10}$ in every case. Proposition~\ref{prop:bifurcation} predicts which solutions they are and how far from the reference state they lie. Writing $d_{\rm ref}$ for the distance from the returned state to the reference equilibrium relative to that equilibrium, Eq.~\eqref{eq:bifurcation_scaling} gives $d_{\rm ref}\propto\rho_{\rm ref}-1$ near the bifurcation point. Figure~\ref{fig:branch_scaling} plots the $1696$ measured distances against $\rho_{\rm ref}-1$, four of them, all on random meshes, outside the axes. Over the $959$ cases with $\rho_{\rm ref}<1.10$ a one-parameter fit gives $d_{\rm ref}=0.202\,(\rho_{\rm ref}-1)$ with $R^{2}=0.96$, and the same slope persists to $\rho_{\rm ref}=1.40$, which covers all but three of the population; those three occupy the highest band of Table~\ref{tab:branch_classification} and do not follow the trend. The three mesh families fall on one line. The spectral radius at the reference equilibrium therefore grades not only the risk of reaching a different solution but, quantitatively, the error that results; its median over the whole population is $1.7\%$.

\paragraph{The alternative states survive mesh refinement}
A separate campaign repeated the same $100$ draws on three refinement levels of every mesh family. Of those draws, $22$ end at a non-reference state on every level of both the Cartesian and the random family, and Table~\ref{tab:branch_refinement} gives the median distance from the intended solution over that common set. Refining by a factor of $63$ in the number of cells leaves it unchanged to within $0.5\%$, where second-order discretisation error would reduce it by a factor of $63$, and changing the mesh topology at fixed resolution moves it no further; only the coarsest random level departs, by $6\%$. A distance between two discrete states that is insensitive to both resolution and topology is not a discretisation artefact, so the alternative states are to be read as a second solution branch of Eq.~\eqref{eq:manufactured_continuous}, of the kind that deflation methods are designed to locate~\cite{Farrell2015DeflationEquations}. The hexagonal family records $4$, $2$, and $32$ non-reference cases on the three levels and so contributes only one draw to the common set; the cause of the discrepancy at its coarser levels is not established here.

\begin{table}[htb]
\centering
\caption{How far the alternative state lies from the intended one, across resolutions and topologies.}
\label{tab:branch_refinement}
\setlength{\tabcolsep}{7pt}
\begin{tabular}{lrr}
\toprule
\rowcolor{gray!12}
Mesh and level & Cells & Median distance \\
\midrule
Cartesian, coarse & $5504$ & $2.1218\%$ \\
Cartesian, medium & $87552$ & $2.1155\%$ \\
Cartesian, fine & $349184$ & $2.1151\%$ \\
\addlinespace[2pt]
Random, coarse & $5504$ & $1.9849\%$ \\
Random, medium & $87552$ & $2.1226\%$ \\
Random, fine & $349184$ & $2.1120\%$ \\
\addlinespace[2pt]
Hexagonal, fine & $403\,259$ & $2.1181\%$ \\
\bottomrule
\end{tabular}
\end{table}

\subsection{Which spectral radius is usable}
\label{subsec:which_rho}

The classification above uses $\rho_{\rm ref}$, which a production calculation does not have. Three evaluation points must be distinguished, and they do not carry the same information. Evaluated at the state the solver returns, the spectral radius certifies nothing: over the $1696$ non-reference states it is below unity in every one, with median $0.9213$ and largest value $0.9964$, as the exchange of stability in Proposition~\ref{prop:bifurcation} requires. This is not accidental: the returned states are locally attracting, whatever solution they are, and the measured radii show it directly. Evaluated at the initial iterate it is both available and informative. The initial-$\rho$ campaign of Section~\ref{subsec:newton_and_mixed} reached the intended solution in all $6000$ cases, but that record does not by itself test the criterion, since every case the rule assigns to Newton succeeds whether or not $\rho_{\rm initial}$ is informative. The test is the amplitude sweep below, over Picard solves alone, where the criterion admitted no counterexample at the amplitude of the main campaigns. Table~\ref{tab:rho_variants} summarises the distinction, which is the operative result of this work: the quantity that decides the outcome is measurable before the solve and useless after it.

\begin{table}[htb]
\centering
\caption{Why the three evaluation points of $\rho$ are not interchangeable, at the amplitude of the main campaigns.}
\label{tab:rho_variants}
\setlength{\tabcolsep}{7pt}
\begin{tabular}{llll}
\toprule
\rowcolor{gray!12}
 & & Available & Separates \\
\rowcolor{gray!12}
Evaluated at & Notation & in practice & the solutions \\
\midrule
initial iterate & $\rho_{\rm initial}$ & yes & yes, \\ &&& within a measured neighbourhood \\
returned state & $\rho_{\rm final}$ & yes & no, below unity in all $1696$ \\
intended solution & $\rho_{\rm ref}$ & no & one-sided, no \\ &&& exception below unity \\
\bottomrule
\end{tabular}
\end{table}

\paragraph{Extent of the neighbourhood}
The predictive value of $\rho_{\rm initial}$ rests on the initial iterate being close to the solution, as it is in the campaigns above. To bound that condition the perturbation amplitude $a$ of Table~\ref{tab:protocol} was swept over four decades at the same mesh sizes and free functions, on a paired subset of $200$ cases per strategy and mesh, the level $10^{-4}$ being the matching cases of the campaigns above. Over the Picard solves that ended without a technical exception, $\rho_{\rm initial}<1$ was followed by a failure to reach the reference equilibrium in none of the cases at amplitude $10^{-4}$, $4$ of $298$ at $10^{-2}$ and $7$ of $255$ at $10^{-1}$, while $\rho_{\rm initial}\geq1$ was followed by recovery in $0$ of $268$ and $20$ of $272$ cases at the last two amplitudes. At $10^{0}$ the radius exceeds unity at every initial iterate, so the criterion predicts failure throughout; $19$ of $440$ solves nonetheless reached the reference, and Newton reached it in only $3.3\%$ of its solves. The criterion is thus exact at the amplitude of the main campaigns, degrades gradually as the initial iterate moves away, and carries no information once it is displaced by an amount comparable to the solution itself. This is the condition the many-query regime supplies: consecutive forward solves of a sampling calculation differ by one move of the sampler, so each begins near the previous solution, and the sweep measures how large a move the prediction tolerates.

\subsection{Newton and the two spectral strategies}
\label{subsec:newton_and_mixed}

\paragraph{Newton-only campaign}
Newton is insensitive to the spectral partition. Table~\ref{tab:final_newton} shows no loss of robustness over the sampled range and a nonlinear work that does not respond to it, including in the bin where Picard fails in every case. At this amplitude Newton reached the intended solution in all $6000$ cases and no non-reference convergence was observed. Every case carries $\rho_{\rm ref}$ in this campaign, because the exceptions of Section~\ref{subsec:numerical_protocol} arose only under Picard, so the $2981$ cases here with $\rho_{\rm ref}\geq1$ are the $2737$ of Table~\ref{tab:branch_classification} together with the $244$ that Picard could not classify.  The dominant variation in cost is instead associated with the mesh. Cartesian meshes are the least expensive, and the higher costs on the hexagonal and random families follow the cost of assembling and solving the Jacobian systems on them: the coloured assembly requires $7$, $10$, and $14$ colours and the Jacobian holds $4.97$, $6.94$, and $6.95$ non-zeros per cell on the Cartesian, hexagonal, and random meshes respectively. One random-mesh case required seven Newton iterations; it is the case discussed after Table~\ref{tab:final_mixed_initial}.

\begin{table}[htb]
\centering
\caption{The same quantities for Newton, which does not respond to the spectral partition.}
\label{tab:final_newton}
\setlength{\tabcolsep}{7pt}
\begin{tabular}{lcc}
\toprule
\rowcolor{gray!12}
Mesh &
$\rho<0.90$ &
$0.90\leq\rho<0.98$ \\
\midrule
Cartesian &
$(1.00,\,3,\,0.0\%)$ &
$(1.19,\,3,\,0.0\%)$ \\
Hexagonal &
$(1.69,\,3,\,0.0\%)$ &
$(1.87,\,3,\,0.0\%)$ \\
Random &
$(2.64,\,4,\,0.0\%)$ &
$(2.80,\,4,\,0.0\%)$ \\

\addlinespace[2pt]
\rowcolor{gray!12}
Mesh &
$0.98\leq\rho<1.02$ &
$\rho\geq1.02$ \\
\midrule
Cartesian &
$(1.09,\,3,\,0.0\%)$ &
$(1.11,\,3,\,0.0\%)$ \\
Hexagonal &
$(1.82,\,3,\,0.0\%)$ &
$(1.83,\,3,\,0.0\%)$ \\
Random &
$(2.83,\,4,\,0.0\%)$ &
$(2.87,\,4,\,0.0\%)$ \\
\bottomrule
\end{tabular}
\end{table}

\paragraph{Initial-$\rho$ mixed campaign}
The initial-$\rho$ strategy reported in Table~\ref{tab:final_mixed_initial} reached the manufactured reference solution for every sampled parameter combination on every mesh. This is consistent with the intended role of the spectral selector: Picard is retained where the initial spectral state indicates a locally contractive regime, and the remaining cases are transferred to Newton. In the critical interval, however, the nonlinear work remains overwhelmingly dominated by Picard. The cases retained by Picard lie close to the contractivity boundary and therefore require substantially more iterations, whereas those transferred to Newton converge in a few steps. The mixed strategy thus removes the reference-solution failures of the Picard-only campaign without removing the high Picard cost of near-critical cases.

The small departure from exclusive Picard use in the low-$\rho_{\rm ref}$ random-mesh regime follows from the distinction between the reference and initial spectral states. The tables are stratified by $\rho_{\rm ref}$, whereas the solver is selected from $\rho_{\rm initial}$. In the single affected case $\rho_{\rm ref}=0.894114$ but $\rho_{\rm initial}=1.076818$, so Newton was selected and converged in seven iterations. The event is consistent with the switching rule and illustrates why the spectral state used for the decision must be distinguished from that of the reference equilibrium.

\begin{table}[htb]
\centering
\caption{What the initial-$\rho$ selector removes, and what it does not.}
\label{tab:final_mixed_initial}
\setlength{\tabcolsep}{7pt}
\begin{tabular}{lcc}
\toprule
\rowcolor{gray!12}
Mesh &
$\rho<0.90$ &
$0.90\leq\rho<0.98$ \\
\midrule
Cartesian &
$(1.22,\,29,\,0.0\%, \,100.0\%)$ &
$(2.44,\,83,\,0.0\%, \,100.0\%)$ \\
Hexagonal &
$(1.96,\,29,\,0.0\%, \,100.0\%)$ &
$(3.83,\,81,\,0.0\%, \,100.0\%)$ \\
Random &
$(2.00,\,30,\,0.0\%, \,99.9\%)$ &
$(3.75,\,86,\,0.0\%, \,100.0\%)$ \\

\addlinespace[2pt]
\rowcolor{gray!12}
Mesh &
$0.98\leq\rho<1.02$ &
$\rho\geq1.02$ \\
\midrule
Cartesian &
$(1.38,\,165,\,0.0\%, \,98.9\%)$ &
$(1.00,\,3,\,0.0\%, \,0.0\%)$ \\
Hexagonal &
$(2.01,\,159,\,0.0\%, \,98.9\%)$ &
$(1.56,\,3,\,0.0\%, \,0.0\%)$ \\
Random &
$(2.81,\,198,\,0.0\%, \,99.0\%)$ &
$(2.09,\,4,\,0.0\%, \,0.0\%)$ \\
\bottomrule
\end{tabular}
\end{table}



\paragraph{Updated-$\rho$ mixed campaign}
Recomputing the spectral information before every nonlinear step preserves the same robustness, as Table~\ref{tab:final_mixed_updated} shows, but does not reduce the nonlinear work relative to the initial-$\rho$ strategy: the mean numbers of nonlinear solves and the Picard participation are essentially unchanged. The computational cost, however, rises sharply, because the eigensolve is repeated along the trajectory. Across the $6000$ WLSQR cases the initial-$\rho$ strategy required $6000$ spectral evaluations and the updated strategy $352\,166$, and one solver switch occurred in the entire updated campaign. Repeated spectral evaluation therefore supplied virtually no additional information for the sampled problems.

\begin{table}[htb]
\centering
\caption{What changes when the spectrum is recomputed at every step, and what does not.}
\label{tab:final_mixed_updated}
\setlength{\tabcolsep}{7pt}
\begin{tabular}{lcc}
\toprule
\rowcolor{gray!12}
Mesh &
$\rho<0.90$ &
$0.90\leq\rho<0.98$ \\
\midrule
Cartesian &
$(7.99,\,29,\,0.0\%, \,100.0\%)$ &
$(22.06,\,83,\,0.0\%, \,100.0\%)$ \\
Hexagonal &
$(13.86,\,29,\,0.0\%, \,100.0\%)$ &
$(36.84,\,81,\,0.0\%, \,100.0\%)$ \\
Random &
$(13.86,\,30,\,0.0\%, \,100.0\%)$ &
$(36.99,\,86,\,0.0\%, \,100.0\%)$ \\

\addlinespace[2pt]
\rowcolor{gray!12}
Mesh &
$0.98\leq\rho<1.02$ &
$\rho\geq1.02$ \\
\midrule
Cartesian &
$(2.04,\,165,\,0.0\%, \,98.9\%)$ &
$(1.00,\,3,\,0.0\%, \,0.0\%)$ \\
Hexagonal &
$(3.07,\,159,\,0.0\%, \,98.9\%)$ &
$(1.80,\,3,\,0.0\%, \,0.0\%)$ \\
Random &
$(17.38,\,198,\,0.0\%, \,99.0\%)$ &
$(2.17,\,4,\,0.0\%, \,0.0\%)$ \\
\bottomrule
\end{tabular}
\end{table}

\subsection{Cost in the many-query regime}
\label{subsec:many_query}

Table~\ref{tab:strategy_comparison} uses the median algorithmic time per case, which excludes the reference and initial spectral diagnostics recorded in every campaign but used only by the spectral strategies, with Newton on the Cartesian mesh ($0.074931\,\mathrm{s}$) as $C_{\rm rel}=1$. The ordering is the same on all three mesh families: Newton, then initial-$\rho$, then Picard, then updated-$\rho$.

Newton has the lowest cost and the smallest nonlinear work everywhere, and both spectral strategies match its reference-solution failure rate. The two perform identical nonlinear work, as they must, since only one of the $6000$ cases changed solver after the first step; they differ only in spectral evaluations, and it is that difference alone that places updated-$\rho$ last. The initial selector pays one evaluation per forward solve, the updated strategy approximately one per nonlinear iteration. The solve counts are not equal units and are not to be compared as such: a Picard solve is one application of the invariant factorisation, whereas a Newton step also requires a coloured Jacobian assembly and a preconditioned Krylov solve, which took between $5$ and $6$ iterations in the main campaigns. Failure rates are evaluated over cases for which $\rho_{\rm ref}$ is available, so the technical exceptions are excluded from both numerator and denominator.

\begin{table}[h]
\centering
\caption{The four strategies on a common cost basis.}
\label{tab:strategy_comparison}
\begin{tabular}{llrrrr}
\toprule
\rowcolor{gray!12}
Mesh & Strategy & $C_{\rm rel}$ & Solves & Spectral & Failure \\
\midrule
Cartesian & Newton          &  1.00 &   3 &  0.0 &  0.0\% \\
          & Initial-$\rho$  &  2.19 &  57 &  1.0 &  0.0\% \\
          & Picard          &  7.47 & 363 &  0.0 & 47.4\% \\
          & Updated-$\rho$  &  8.42 &  57 & 56.7 &  0.0\% \\
\midrule
Hexagonal & Newton          &  1.64 &   3 &  0.0 &  0.0\% \\
          & Initial-$\rho$  &  3.38 &  55 &  1.0 &  0.0\% \\
          & Picard          &  8.90 & 361 &  0.0 & 48.0\% \\
          & Updated-$\rho$  & 12.85 &  55 & 55.1 &  0.0\% \\
\midrule
Random    & Newton          &  2.55 &   4 &  0.0 &  0.0\% \\
          & Initial-$\rho$  &  4.11 &  64 &  1.0 &  0.0\% \\
          & Picard          &  9.77 & 326 &  0.0 & 44.8\% \\
          & Updated-$\rho$  & 32.98 &  64 & 64.4 &  0.0\% \\
\bottomrule
\end{tabular}
\end{table}

\paragraph{Projection to the many-query regime}
The intended application is expected to require between $10^{4}$ and $10^{6}$ forward solutions. To show how the measured per-case differences accumulate at those counts, Table~\ref{tab:many_query_projection} gives a linear serial projection based on the median algorithmic time per attempted WLSQR case. The projection holds the present discretisation fixed and varies only the number of forward solves; it is therefore not an extrapolation to a larger mesh. For Picard the projected counts are attempted forward solves and do not imply that all attempts reach the manufactured reference solution. The second row of each block is Newton with one spectral evaluation per forward solve, the configuration in which the radius is used to anticipate the outcome rather than to choose the solver: its per-case time is the median over the Newton campaign of the sum of the nonlinear solve and the spectral evaluation recorded in that same campaign, so it is measured on the same cases and not composed from separate runs.

One spectral evaluation costs $0.0580$, $0.0844$, and $0.0936\,\mathrm{s}$ on the Cartesian, hexagonal, and random meshes, which is $0.77$, $0.69$, and $0.49$ of a complete Newton solve there. Anticipating the outcome of every forward solve raises the cost by a factor of $1.80$, $1.73$, and $1.49$, each the median of the per-case sum rather than the sum of the two medians. Using the same information to select a solver rather than to anticipate the outcome costs more, because the subcritical cases it assigns to Picard require many more nonlinear solves than Newton. Recomputing the spectrum along the trajectory costs about $4$ to $8$ times the cost of the initial selector, and returns the single altered decision already noted.

The measurements separate what the spectral radius predicts from where it can be evaluated. Residual convergence does not identify which discrete solution a semilinear iteration has reached; the alternative states are genuine solutions that survive a $63$-fold refinement and lie where the bifurcation analysis places them; and the spectral radius at the returned state does not distinguish them. Evaluated at the initial iterate it does, without exception at the amplitude of the main campaigns, and the amplitude sweep measures the neighbourhood in which that holds. The evaluation costs less than one forward solve, which is affordable in the many-query regime precisely because it reuses the factorisation that the regime already amortises. Two limits of this evidence should be stated. At the problem size and initial amplitude of the main campaigns Newton is the least expensive strategy and reaches the intended solution in every case, so the spectral radius earns its place as a predictor of the outcome rather than as a rule for choosing the solver; and whether reuse of the invariant Picard operator alters that ranking at substantially larger problem sizes cannot be inferred from a projection at fixed problem size.


\begin{table}[htb]
\centering
\caption{How the per-case differences accumulate over a many-query calculation.}
\label{tab:many_query_projection}
\setlength{\tabcolsep}{5pt}
\begin{tabular}{llrrrr}
\toprule
\rowcolor{gray!12}
 & & Median time & & & \\
\rowcolor{gray!12}
Mesh & Strategy & per case & $10^4$ solves & $10^5$ solves & $10^6$ solves \\
\midrule
Cartesian & Newton & $0.075\,\mathrm{s}$ & $0.21\,\mathrm{h}$ & $2.08\,\mathrm{h}$ & $0.87\,\mathrm{d}$ \\
 & Newton $+$ one $\rho$ & $0.135\,\mathrm{s}$ & $0.37\,\mathrm{h}$ & $3.75\,\mathrm{h}$ & $1.56\,\mathrm{d}$ \\
 & Initial-$\rho$ & $0.164\,\mathrm{s}$ & $0.46\,\mathrm{h}$ & $4.57\,\mathrm{h}$ & $1.90\,\mathrm{d}$ \\
 & Picard & $0.560\,\mathrm{s}$ & $1.56\,\mathrm{h}$ & $15.55\,\mathrm{h}$ & $6.48\,\mathrm{d}$ \\
 & Updated-$\rho$ & $0.631\,\mathrm{s}$ & $1.75\,\mathrm{h}$ & $17.52\,\mathrm{h}$ & $7.30\,\mathrm{d}$ \\
\addlinespace[2pt]
Hexagonal & Newton & $0.123\,\mathrm{s}$ & $0.34\,\mathrm{h}$ & $3.41\,\mathrm{h}$ & $1.42\,\mathrm{d}$ \\
 & Newton $+$ one $\rho$ & $0.212\,\mathrm{s}$ & $0.59\,\mathrm{h}$ & $5.90\,\mathrm{h}$ & $2.46\,\mathrm{d}$ \\
 & Initial-$\rho$ & $0.253\,\mathrm{s}$ & $0.70\,\mathrm{h}$ & $7.04\,\mathrm{h}$ & $2.93\,\mathrm{d}$ \\
 & Picard & $0.667\,\mathrm{s}$ & $1.85\,\mathrm{h}$ & $18.53\,\mathrm{h}$ & $7.72\,\mathrm{d}$ \\
 & Updated-$\rho$ & $0.963\,\mathrm{s}$ & $2.68\,\mathrm{h}$ & $26.75\,\mathrm{h}$ & $11.15\,\mathrm{d}$ \\
\addlinespace[2pt]
Random & Newton & $0.191\,\mathrm{s}$ & $0.53\,\mathrm{h}$ & $5.32\,\mathrm{h}$ & $2.21\,\mathrm{d}$ \\
 & Newton $+$ one $\rho$ & $0.286\,\mathrm{s}$ & $0.79\,\mathrm{h}$ & $7.94\,\mathrm{h}$ & $3.31\,\mathrm{d}$ \\
 & Initial-$\rho$ & $0.308\,\mathrm{s}$ & $0.85\,\mathrm{h}$ & $8.55\,\mathrm{h}$ & $3.56\,\mathrm{d}$ \\
 & Picard & $0.732\,\mathrm{s}$ & $2.03\,\mathrm{h}$ & $20.34\,\mathrm{h}$ & $8.47\,\mathrm{d}$ \\
 & Updated-$\rho$ & $2.471\,\mathrm{s}$ & $6.87\,\mathrm{h}$ & $68.65\,\mathrm{h}$ & $28.60\,\mathrm{d}$ \\
\bottomrule
\end{tabular}
\end{table}

